\documentclass{report}
\usepackage{amsmath,amssymb}
\setlength{\parindent}{0mm}
\setlength{\parskip}{1em}
\begin{document}
\begin{center}
\rule{6in}{1pt} \
{\large hassane sadok \\
{\bf A new approach to Conjugate Gradient convergence}}

L M P A \\ Universit� du Littoral \\ 50 rue F Buisson BP699 \\ F-62228 Calais Cedex \\ France
\\
{\tt sadok@lmpa.univ-littoral.fr}\\
Mohammed Bellalij\end{center}

Known as one of the best iterative methods for solving symmetric
positive definite linear systems, CG generates as FOM an
Hessenberg matrix which is symmetric then triangular.

This
specific structure may be really helpful to understand how does
behave the convergence of the conjugate gradient method and its
study gives an interesting alternative to Chebyshev polynomials.
The talk deals about some new bounds on residual norms and error
$A$-norms using essentially the condition number.

We will show how to derive a bound of the $A$- norm of the error by
solving a constrained optimization problem using Lagrange multipliers.


\end{document}
